\section{问题从哪来：Atiyah 的两问与 Jones 多项式的缺口}

\subsection{一个 1988 年的追问}

本文开篇引的是 Atiyah 在“去年 Hermann Weyl 讨论会”上提的两个问题 \src{351}。原话只能从 Witten 的转述读到（他说 Atiyah “might roughly be described” 这样问），原文文献 [1] 本身不在引用范围内可查。两问是：

\begin{enumerate}[leftmargin=2em]
\item 能否给 Donaldson 的四维流形理论一个物理诠释？
\item 能否给 Jones 多项式一个内蕴的三维定义？
\end{enumerate}

第一问在同一年由 Witten 自己的另一篇文章解决——用扭曲的 $N=2$ 超对称 Yang--Mills 理论 \src{351}，即参考文献 [4]（\emph{Topological quantum field theory}, CMP 117 (1988) 353）。本文处理第二问。

\begin{keypoint}
\textbf{注意}：中译本每节开头的导语框里有一句“本文解决第二个问题”。这句话是\textbf{译者所加}，不在原刊引言里。原文在引言与摘要中只说这篇理论 “gives a natural framework” 给 Jones 多项式 \src{352}，用词比“解决”弱。引用时必须回到正文。
\end{keypoint}

\subsection{Jones 多项式缺什么}

Jones 在 1983--1987 年间给出了一个纽结不变量，它有几种等价的、非常漂亮的\textbf{二维}定义：穿过投影图的切片（skein）关系、辫群表示、von Neumann 代数迹 \src{352}（文献 [5]）。但它没有\textbf{内蕴三维}的定义——不知道它在三维空间里“是什么”。

Witten 用的类比很精确 \src{352}：1930 年代人们知道量子电动力学是相对论性的，却写不出一个明显相对论性的表述；直到有了量子场论的语言，这件事才“显然”。Jones 多项式当时就处在这个位置：二维上算得出来，但不知道三维里的对象是什么。

\subsection{物理侧的动机：二维来源太多}

同期物理学家发现纽结多项式从好几个地方冒出来 \src{352}：
\begin{itemize}[leftmargin=1.6em]
\item 可解格点模型（统计力学，文献 [7] Kauffman）；
\item Yang--Baxter 方程与量子群（文献 [8] Turaev，[35] Drinfeld）；
\item 二维共形场论的单值性（文献 [11] Tsuchiya--Kanie）。
\end{itemize}

“来源太多”本身是个问题：它们之间是什么关系？Witten 的原话是，要“为这种多样性带来秩序、找到统一主题” \src{352}。更具体地，他问：\textbf{二维共形场论里，究竟是什么东西是三维的？}这个问题后来被证明是整篇文章的种子：三维的 Chern--Simons 理论沿曲面切开给出的 Hilbert 空间，正是二维共形场论的共形块。

\subsection{“拓扑不变性 = 广义协变”}

引言里有一个容易被跳过的观念转折 \src{353}：通常说一个理论“拓扑不变”，指的是配分函数不依赖度规；而广义相对论里“对度规积分”只是实现这种不变性的一种方式。Chern--Simons 理论提供了另一种：\textbf{作用量里根本没有度规}，所以不需要对度规积分，量子化出来的量自动不依赖度规。

这一步是整篇文章的立足点：如果作用量含度规，量子化后一般会得到依赖度规的量；不含度规，才有资格谈拓扑不变量。

\subsection{本文的日期与身份}

文末记：1988 年 9 月 12 日收稿，9 月 27 日修改；由 A. Jaffe 送审。p.351 的脚注标明本文是 1988 年 7 月 Swansea 国际数学物理大会报告的扩充版。所以它的“学术日历”是：先有报告，再有这篇论文；文中的一些说法（例如“量子群”评论）带着报告式的前瞻语气，不宜当定理读。

\subsection{摘要里的四个动词}

摘要 \src{351} 说：2+1 维量子 Yang--Mills 在作用量纯为 Chern--Simons 项时可以精确求解（“can be exactly soluble”）；这给 Jones 多项式一个自然框架（“gives”）；Jones 多项式可以推广到任意三维流形（“can be generalized”）；并且是可由手术计算的（“computable”）；还给 1+1 维共形场论带来“出人意料的新启示”（“shed new light”）。四个动词里没有 “prove”。这提示了全文的性质：\textbf{构造 + 论证 + 计算方案}，不是一套公理化的证明。

\begin{boundary}
\begin{itemize}[leftmargin=1.4em]
\item 已确证：Atiyah 两问的转述文本、第一问由 [4] 解决、Jones 缺三维定义、二维来源的三条线索、Swansea 版本、收稿日期。
\item 我们补的：把“二维来源过多”解释为“统一动机”，原文只说“带来秩序”。
\item 未证/待核：Atiyah 原报告的会议年份与文集信息（原文只写“去年”）；“本文解决第二问”一语不见于原刊。
\end{itemize}
\end{boundary}
